%% ma: L12-B04-0006
%% lop: 12
%% chuong: 1
%% bai: 4
%% dang: 12.4.6
%% loai: DS
%% muc: [TH, VD, TH, VDC]
%% dap_an: "SĐĐS"
%% nguon: "Ảnh giáo viên gửi 11/10/2026 (Đề kiểm tra giữa HK1 Toán 12, THPT Việt Đức 2026-2027, mã đề 0103), Phần II Câu 4"
%% kien_thuc: "Bài 4 (tâm đối xứng, tiệm cận xiên của hàm phân thức bậc hai trên bậc nhất), Bài 3 (tiệm cận xiên), Bài 8 (vectơ, tọa độ), khoảng cách từ điểm đến đường thẳng (lớp 10)"
%% trang_thai: cho-duyet
%% ly_do: "Dạng mới đề xuất, chờ giáo viên duyệt."
%% ghi_chu: "Đáp án tính được S Đ Đ S. Ý d dùng tính chất: hai dây cung cùng nhận một điểm làm trung điểm thì song song, nên hình bình hành nội tiếp hypebol có tâm là tâm đối xứng của hypebol (câu khó, mức VDC)."
\begin{ex}
Cho hàm số $y=\dfrac{-x^2+x-2}{x+2}$ có đồ thị $(C)$. Đường thẳng $\Delta\colon x+2y=-2$ cắt $(C)$ tại hai điểm phân biệt $A$, $B$.
\choiceTF
{Tâm đối xứng của $(C)$ là điểm $I(-2;-5)$.}
{\True Diện tích tam giác $OAB$ bằng $3$.}
{\True Đường tiệm cận xiên của $(C)$ là đường thẳng $y=-x+3$.}
{Gọi $C$, $D$ là hai điểm thuộc đồ thị $(C)$ sao cho $ABCD$ là hình bình hành. Khoảng cách giữa hai đường thẳng $AB$ và $CD$ là $5$.}
\loigiai{Ta có $y=-x+3-\dfrac8{x+2}$.

a) Tiệm cận đứng $x=-2$, tiệm cận xiên $y=-x+3$ nên tâm đối xứng $I(-2;5)$, không phải $(-2;-5)$. Sai.

c) $\lim\limits_{x\to\infty}\left(y-(-x+3)\right)=\lim\limits_{x\to\infty}\dfrac{-8}{x+2}=0$ nên $y=-x+3$ là tiệm cận xiên. Đúng.

b) Phương trình hoành độ giao điểm: $\dfrac{-x^2+x-2}{x+2}=-\dfrac{x+2}2$ ($x\ne-2$) $\Leftrightarrow-x^2+6x=0\Leftrightarrow x=0$ hoặc $x=6$. Vậy $A(0;-1)$, $B(6;-4)$ và $S_{OAB}=\dfrac12|0\cdot(-4)-6\cdot(-1)|=3$. Đúng.

d) Hình bình hành nội tiếp một hypebol phải nhận tâm đối xứng $I(-2;5)$ của hypebol làm tâm, nên $C=2I-A=(-4;11)$, $D=2I-B=(-10;14)$ (kiểm tra thuộc $(C)$). Khoảng cách giữa $AB$ và $CD$ bằng hai lần khoảng cách từ $I$ đến $AB\colon x+2y+2=0$, tức $2\cdot\dfrac{|-2+10+2|}{\sqrt5}=4\sqrt5\ne5$. Sai.}
\end{ex}
